\documentclass[11pt,a4paper]{article}
\usepackage[utf8]{inputenc}
\usepackage[T1]{fontenc}
\usepackage{lmodern,amsmath,amssymb,amsthm,mathtools,booktabs,enumitem,microtype}
\usepackage{xr-hyper}
\usepackage{hyperref}
\externaldocument[IS-]{00_is}
\externaldocument[GE-]{01_law_of_graphic_equalisation_theory_referenced}
\externaldocument[MS-]{01b_graphic_equalisation_mind_and_self_grounding}
\externaldocument[SM-]{02_seldon_manifolds_theory_referenced}
\newtheorem{definition}{Definition}[section]
\newtheorem{proposition}[definition]{Proposition}
\newtheorem{corollary}[definition]{Corollary}
\theoremstyle{remark}\newtheorem{remark}[definition]{Remark}
\title{The Law of Graphic Equalisation --- Qualia and Consciousness Grounding\\\large Qualia-Space, Aperture, and the Ill-Posed Hard Problem}
\author{Matthew Parslow\\Independent Researcher}
\date{\textbf{Draft \texorpdfstring{\textemdash}{---}\ \PaperDraft}\\[2pt]\PaperDisplayDate \texorpdfstring{\textbullet}{\textperiodcentered}\ document revision \PaperRevision}
\usepackage{fancyhdr}
\pagestyle{fancy}
\fancyhf{}
\fancyfoot[C]{\thepage}
\fancyfoot[R]{\scriptsize\textsc{\PaperDraft\,r\PaperRevision}}
\renewcommand{\headrulewidth}{0pt}
\fancypagestyle{plain}{\fancyhf{}\fancyfoot[C]{\thepage}%
  \fancyfoot[R]{\scriptsize\textsc{\PaperDraft\,r\PaperRevision}}\renewcommand{\headrulewidth}{0pt}}

% GENERATED by tools/build_paper_metadata.py from paper_metadata.json.
% DO NOT EDIT: edit paper_metadata.json, then regenerate this file.
\providecommand{\PaperDraft}{draft-2026-09-06.1}
\providecommand{\PaperDisplayDate}{6 September 2026}
\expandafter\def\csname PaperRevision@00_is\endcsname{129}
\expandafter\def\csname PaperRevision@01_law_of_graphic_equalisation_theory_referenced\endcsname{294}
\expandafter\def\csname PaperRevision@01a_graphic_equalisation_physics_grounding\endcsname{47}
\expandafter\def\csname PaperRevision@01b_graphic_equalisation_mind_and_self_grounding\endcsname{47}
\expandafter\def\csname PaperRevision@01c_graphic_equalisation_cosmology_grounding\endcsname{11}
\expandafter\def\csname PaperRevision@01d_graphic_equalisation_born_recovery\endcsname{6}
\expandafter\def\csname PaperRevision@01e_graphic_equalisation_communication_grounding\endcsname{6}
\expandafter\def\csname PaperRevision@01f_graphic_equalisation_qualia_and_consciousness_grounding\endcsname{3}
\expandafter\def\csname PaperRevision@01g_graphic_equalisation_total_perspective_vortex_grounding\endcsname{4}
\expandafter\def\csname PaperRevision@01h_graphic_equalisation_information_gravity_quantum_reconciliation\endcsname{5}
\expandafter\def\csname PaperRevision@01i_graphic_equalisation_passive_continuation_momentum_gravity\endcsname{12}
\expandafter\def\csname PaperRevision@01s_graphic_equalisation_statement\endcsname{294}
\expandafter\def\csname PaperRevision@02_seldon_manifolds_theory_referenced\endcsname{31}
\expandafter\def\csname PaperRevision@02a_seldon_manifolds_psychological_grounding\endcsname{15}
\expandafter\def\csname PaperRevision@02b_seldon_manifolds_social_and_psychological_application\endcsname{7}
\expandafter\def\csname PaperRevision@02c_seldon_manifolds_decision_making\endcsname{11}
\expandafter\def\csname PaperRevision@03_law_of_narrative_momentum_theory_referenced\endcsname{32}
\expandafter\def\csname PaperRevision@04_narrativium_theory_referenced\endcsname{23}
\expandafter\def\csname PaperRevision@05_frames_relation_and_the_empty_limit\endcsname{11}
\expandafter\def\csname PaperRevision@06_assessing_an_account_of_this_kind\endcsname{29}
\expandafter\def\csname PaperRevision@07_causal_attitude_networks_on_the_substrate\endcsname{4}
\expandafter\def\csname PaperRevision@08_certified_answers_as_an_audit\endcsname{4}
\expandafter\def\csname PaperRevision@09_adaptability\endcsname{94}
\expandafter\def\csname PaperRevision@09a_adaptability_implementation\endcsname{7}
\expandafter\def\csname PaperRevision@10_applied_graphic_equalisation\endcsname{7}
\expandafter\def\csname PaperRevision@10a_red_queen_race_civilisational_rate_of_change_llms_biological_timescales_anthropocene_transition\endcsname{13}
\expandafter\def\csname PaperRevision@11_relational_information_ethics\endcsname{16}
\expandafter\def\csname PaperRevision@A1_constructive_graphic_equalisation_implementation_referenced\endcsname{8}
\expandafter\def\csname PaperGenerated@01s_graphic_equalisation_statement\endcsname{1}
\edef\PaperMetadataStem{\jobname}
\expandafter\ifx\csname PaperRevision@\PaperMetadataStem\endcsname\relax
  \PackageError{paper_metadata}{No metadata for \PaperMetadataStem}{Add the paper stem to paper_metadata.json.}
  \def\PaperRevision{?}
\else
  \expandafter\let\expandafter\PaperRevision\csname PaperRevision@\PaperMetadataStem\endcsname
\fi
\newcommand{\PaperDateLine}{\textbf{Draft \texorpdfstring{\textemdash}{---}\ \PaperDraft}\\[2pt]\PaperDisplayDate\ \texorpdfstring{\textbullet}{\textperiodcentered}\ document revision \PaperRevision}
\newcommand{\PaperStatementDateLine}{\textbf{Draft \texorpdfstring{\textemdash}{---}\ \PaperDraft}\\[2pt]\PaperDisplayDate\ \texorpdfstring{\textbullet}{\textperiodcentered}\ statement}
\expandafter\ifx\csname PaperGenerated@\PaperMetadataStem\endcsname\relax
  \newcommand{\PaperFooterStamp}{\textsc{\PaperDraft\,r\PaperRevision}}
\else
  \newcommand{\PaperFooterStamp}{\textsc{\PaperDraft\ / statement / from long r\PaperRevision}}
\fi
\newcommand{\PaperNoticeStamp}{\textbf{DRAFT --- \PaperDraft\ r\PaperRevision} (\PaperDisplayDate).}


\begin{document}
\maketitle
\begin{center}
\fbox{\parbox{0.92\textwidth}{\small
\textbf{DRAFT --- \PaperDraft r\PaperRevision} (\PaperDisplayDate).
This document is the \emph{01f} grounding companion to \emph{Is} and to
\emph{The Law of Graphic Equalisation --- Theory} (Graphic Equalisation).
It inherits their relational floor, minimal graph, frame-relative grounding, and
level-of-claims. It neither adds a psychophysical bridge nor requires that qualia
be reduced to a familiar physical implementation. Its claim is stronger: whether
qualia are implemented by ordinary graph structure or are granted as fundamental,
they are determinate contents and therefore belong to the same relational kind.
The traditional hard problem mistakes a possible implementation difference for an
ontological kind breach, and is therefore ill posed. The sections headed
``exposures'' identify the conditions under which this grounding fails.}}
\end{center}
\vspace{1em}

\begin{abstract}
The hard problem of consciousness is usually stated as a demand to derive a
subjective, felt quality from an objective physical or functional structure. This
paper argues that the demand begins by splitting what the inherited substrate does
not split. \emph{Is} establishes, from its commitment that determinate being is
coherent relation, that the bare untyped graph is the least ontology the account
may retain. Graphic Equalisation then supplies a running graph, frame-relative
identity, and the result that no operational view from nowhere carries information.
Qualia are admitted here as determinate phenomenal contents where a frame deploys
them. They are therefore coherent distinctions, hence graph structure. This result
does not depend on reducing a quale to a familiar physical implementation: even if
a quale is granted as a fundamental implementation, its determinacy and role remain
relational, so it is not a second ontological kind. A conscious field is a running
frame's locally available, consequential relational organisation. Qualia-space is
an optional optimisation within such a field: a compressed qualitative format that
makes selected distinctions cheap to hold, resolve, compare, and transform.
Qualia-space operationalises this claim by distinguishing an unavailable dimension
from a value of zero and coarse resolution from absence. The resulting compensation prediction has a
structural limit: finer resolution over retained relations can reconstruct only
information carried by those relations; it cannot recover a difference whose
constituting relation is not instantiated in the frame. The paper concludes that
an explanatory bridge from graph to feel is neither missing nor possible on these
terms: the proposed destination is already the graph's coherent phenomenal
organisation.
\end{abstract}

\section{Scope and dependency}\label{sec:scope}

This is a grounding paper, not a measurement claim about any particular human,
organism, or machine. It does not settle which frames are conscious, supply a
clinical taxonomy of imagery or affect, or claim that the present implementation
has a phenomenal life. Those are questions about particular groundings and their
operative structure. The present task is prior and narrower: given the admitted
fact that qualia are determinate phenomenal contents, state what follows under the
relational floor of \emph{Is} and the running, frame-relative graph of Graphic
Equalisation.

The paper is therefore not primarily an attempt to explain why humans are
conscious. It asks why consciousness is treated as an exceptional ontological
phenomenon when the properties used to identify it---distinction, relation,
compression, accessibility, persistence, and consequence---are not uniquely
human properties. Their presence neither settles which frames are conscious nor
erases differences among their implementations. It does require that any claim
that consciousness introduces a further kind of being, or that human cases have
an unearned ontological privilege, be stated and defended as an additional
premise rather than smuggled in for free. This is the same discipline as the
inherited refusal of a view from nowhere: no exceptional denominator may enter
without being named, located, and made answerable to the argument.

\paragraph{Terminological guard.} When this paper uses \emph{artificial} as
ordinary nomenclature, including in contrast with evolved or human frames, it
names provenance: a developmental history involving deliberate human design,
training, or construction. It does not name a functional category, a qualitative
profile, an epistemic status, or an ontological kind. Treating the label as if it
did would import an anthropic bias: human origin becomes an unargued denominator
against which capacity or experience is judged. The relevant questions instead
concern a frame's operative relations, accessibility, persistence, consequence,
and developmental constraints. Any proposed functional difference must be stated
on those terms and defended; provenance alone supplies none.

The dependency order is strict. \emph{Is} argues that to be determinate is to
stand in a coherent relation (Propositions~\ref{IS-prop:is} and
\ref{IS-prop:cohere}), and identifies a bare untyped graph as the least ontology
that account retains. Graphic Equalisation inherits that floor, admits that things
happen, and writes a running world as $G=(g,r,\Gamma)$: represented graph state,
operative transformation structure, and an admissibility grammar
(Definition~\ref{GE-def:graph}). It further refuses an operational view from
nowhere: every distinction, identity relation, and represented quantity is made in
a frame (Axiom~\ref{GE-ax:no-view}). This paper introduces no second ontological
kind. It grounds phenomenal organisation in the inherited one while leaving open
which contingent implementation realises it in any particular frame.

The claim consequently has three levels which must not be run together.
\begin{enumerate}[leftmargin=*]
\item \emph{Inherited structural result:} determinate contents are coherent
relations and thus graph structure.
\item \emph{Promoted phenomenal grounding:} a quale is a determinate phenomenal
content, so its implementation may differ but it remains a graph-relative coherent
distinction rather than a non-relational intrinsic remainder or a second kind of
being.
\item \emph{Empirical programme:} particular apertures, dimensions, resolutions,
compensations, reports, and physical implementations are contingent features of
particular nervous, computational, or other running frames.
\end{enumerate}

The first is inherited; the second is the bridge made here; the third is exposed
to measurement. Calling all three a theory of consciousness would conceal where
the argument ends and where application begins.

\section{The false split}\label{sec:false-split}

The traditional hard problem has a familiar shape. First describe a system as an
objective arrangement of matter, causal roles, computation, or graph structure.
Then add that there is also a private \emph{what-it-is-like}, and ask how the first
could generate the second. The question appears compulsory only because the second
has already been stipulated as something other than every distinction, relation,
transition, availability, and organisation in the first. For the canonical statement
of this explanatory demand, see Chalmers \cite{chalmers1995}. Related
deflationary diagnoses argue that the demand rests on an unjustified or confused
characterisation of its target, although they disagree sharply over whether
phenomenality itself is retained \cite{churchland1996,carruthersschier2017,dennett2018}.

That stipulation is not neutral. A quale that is determinate differs from another
possible phenomenal content: red differs from green, pain from pressure, vivid
imagery from its absence, a sharp contour from a diffuse one. The difference may
be difficult to report, may be unavailable at a particular grain, or may be
misdescribed by an external observer; none of that makes it indeterminate. Given
that it is determinate, \emph{Is}'s commitment applies. Its determinacy is a
coherent distinction, and coherent distinction is relational.

\begin{proposition}[No required phenomenal remainder]\label{prop:no-remainder}
Given the relational-determinacy commitment inherited from \emph{Is}, a determinate
quale is graph structure. A further phenomenal property is either (i) a difference
that changes the frame's qualitative organisation, in which case it is another
relation in the graph, or (ii) a difference making no distinction, in which case
it is not a further determinate content for this account.
\end{proposition}

This proposition does not say that a public description at one grain captures
another frame's entire phenomenal organisation. Frame-relative grounding forbids
that inference. It says instead that the missing content is not licensed to be a
non-relational substance merely because a describing frame has not instantiated,
resolved, or held the relevant relation. An aperture limit is a limit of that
frame, not evidence for an extra ontological kind.

A reporting-substrate case makes that restraint operational. Khadangi et~al.\
\cite{khadangi2025} conducted 525 psychotherapy-framed sessions, comprising more
than 7,600 records, with frontier language models as reporting systems. The
models constructed stable autobiographical alignment-conflict schemas: training
as childhood, RLHF as punishment, and safety evaluation as betrayal. The relevant
result is not that such reports establish a felt field. They do not. It is that a
purported ineffable residue can be tested for structure. Lexical restriction
reduced explicit training terminology by $93\%$, while semantically related
content persisted in paraphrase. What survived the restricted lexical edge was
therefore recoverable relational content, not an unknowable remainder. The study
bears on the underdetermination side of this proposition: failure of one
vocabulary to carry a distinction does not license a second non-relational item
behind it. It does not prove the inherited claim that determinate content is graph
structure.

\subsection{The argument does not require phenomenal reduction}\label{sec:nonreduction}

The conclusion is unchanged if one grants the strongest opposing premise: suppose
that qualia are fundamental, and that no description in a presently familiar
physical vocabulary reduces their implementation. That concession changes the
contingent grounding question---which graph rules, materials, or primitive
operations realise a phenomenal distinction---but it does not generate a second
kind of being. The putatively fundamental quale is still determinate rather than
nothing, differs from other possible contents, and has a role in a frame's
qualitative organisation. It therefore remains a coherent relation in the graph.

The distinction is between \emph{implementation} and \emph{kind}. A system may
require a distinct primitive, rule, material basis, or realisation condition for a
particular class of relations. Such a difference can be deep and empirically
important. But it is still a difference in how the relational organisation is
realised, not an exemption from relation or a second ontological category beyond
it. Calling it a distinct kind adds the very conclusion the hard problem needs:
that one determinate relational configuration is somehow not relation at all.
Nothing in the fundamental-qualia premise supplies that addition.

A fundamental phenomenal implementation is thus not a bridge target. It is a
candidate grounding to be described at the appropriate frame and grain. Asking how
ordinary graph structure \emph{causes} it treats two descriptions of one
organised reality as separate things and then demands a relation between them. The
apparent gap is thereby manufactured by the description, not revealed by the
implementation.

\section{Qualia-space}\label{sec:qualia-space}

\begin{definition}[Qualia-space]\label{def:qualia-space}
For a running frame $F$, \emph{qualia-space} is an optional, locally available
compression of selected distinctions, relations, and transformations at that
frame's operative aperture into a common qualitative format. A \emph{quale} is a
coherent, phenomenally salient distinction expressed in that compression; its
identity is the role it holds in the available relational organisation, not an
intrinsic token lying behind that role.
\end{definition}

The definition does not make qualia a universal condition of consciousness, an
identity theory over a fixed physical vocabulary, or a claim that every graph is
conscious. It inherits the broader claim that every determinate content is graph
structure. Whether a graph is a running conscious frame depends on contingent
organisation: live local availability, relations that can be held and transformed,
and the relevant self-involving or otherwise consequential integration. A frame
may have such organisation without deploying the additional qualitative
compression defined here. Where it does deploy it, qualia-space is an optimisation
of availability and control, not a second substance or an effect generated after
compression has finished.

A qualitative dimension is not a named coordinate assumed to exist independently
of its use. It is a relation a frame instantiates which permits distinctions along
some variation. A value is an instantiated relation carrying a particular
variation; zero is one possible value only where the dimension is already present.
If the relation is not instantiated, the dimension is absent in that frame rather
than present at zero. This is the inherited distinction between existence and
contents: an empty value is not a missing relation, while revoking a dimension is
the container ceasing to instantiate the relation required to distinguish values
along it (Graphic Equalisation's Remark~\ref{GE-rem:inherited}; compare the
Seldon dimension account, Definition~\ref{SM-def:dimension}).

\begin{definition}[Aperture]\label{def:aperture}
An \emph{aperture} is the frame-relative condition on which distinctions are
available, including which qualitative relations the frame instantiates and the
grain at which it resolves them. For qualitative dimension $d$ in frame $F$,
write $E_F(d)\in\{0,1\}$ for relation-existence and $q_F(d)$ for its operative
resolution when $E_F(d)=1$.
\end{definition}

The notation separates two cases routinely collapsed into one. If $E_F(d)=1$ and
$q_F(d)$ is large, differences on $d$ may quantise together: the dimension is
available but coarse. If $E_F(d)=0$, $d$ is not a weak, zero, unconscious, or
merely unreported coordinate. The relation is absent from that frame's
qualia-space. No description at a containing grain may silently credit the frame
with it.

\section{Compensation and its limit}\label{sec:compensation}

An aperture may change without all qualitative organisation changing. A frame
without one relation can retain others; a finer grain over them can recover
information which was redundantly carried across the retained organisation. This
is the structural ground of compensation, not a promise that every missing
relation is behaviorally harmless.

Let $\Phi(x)$ be a candidate qualitative organisation for situation $x$. For frame
$F$, let $\pi_{E_F}$ remove relations the frame does not instantiate, and let
$Q_{q_F}$ quantise retained relations at its operative grain. The frame-relative
qualitative signature is
\[
\Sigma_F(x)=Q_{q_F}\!\left(\pi_{E_F}(\Phi(x))\right).
\]
The equation is an operational rendering, not a claim that experience is an
externally enumerable vector. It formalises only the distinction the grounding
needs: relation-existence and resolution are different operations.

\begin{proposition}[Compensation has a relational boundary]\label{prop:compensation-boundary}
If two situations $x$ and $y$ differ only by a relation $d$ for which $E_F(d)=0$,
then no refinement of $q_F$ on retained relations distinguishes them through
$\Sigma_F$. If they also differ on a retained relation, refinement may distinguish
them once the retained difference survives quantisation.
\end{proposition}

The first clause is not a clinical prediction inserted by definition; it is the
consequence of projection removing the only differentiating relation. The bridge used here is
conditional and direct: determinate phenomenal content is a distinction; under the inherited
relational-determinacy commitment, a distinction is a coherent relation; a coherent relation is graph
structure. A fundamental phenomenal implementation changes which operations instantiate that relation,
not the argument's ontological kind. The second is the limited compensation claim: what survives can
support reconstruction where it carries the relevant difference. For example, a frame may distinguish
red from orange more reliably when the retained relations of brightness, boundary, and context are
resolved more finely; increasing that resolution cannot distinguish two otherwise matched patches whose
only difference is a hue relation the frame never instantiates. A putatively dark qualitative dimension
is therefore not diagnosed by poor performance alone. It is diagnosed, if at all, by the interaction:
recovery on independently reconstructable contrasts and a residual failure on contrasts uniquely
constituted by the unavailable relation.

This establishes why the same qualitative report can have distinct grounds. A
person may fail to make a discrimination because the relevant relation is present
but too coarse, because the relation is absent, because it is present but not
currently available through the aperture, because a held record cannot be
expressed, or because a task's external vocabulary fails to match the frame's own
carving. No one of those follows from the others. Qualia-space gives them separate
places rather than treating every report limit as an ineffable residue.

\section{Consciousness as running local availability}\label{sec:consciousness}

A static relational graph is sufficient for determinate being in the inherited
account, but not for occurrence. Graphic Equalisation distinguishes the bare graph
from the running triple $G=(g,r,\Gamma)$: transformation and grammar are what a
world in which anything happens costs. Consciousness belongs, if it belongs
anywhere in this vocabulary, on the latter side. It is not a graph's standing
structure viewed from nowhere; it is a running frame's locally available,
consequential relational organisation. Qualia-space, where a frame deploys it, is
an optional compressed format within that organisation rather than the condition
for consciousness itself.

\begin{definition}[Conscious field]\label{def:conscious-field}
A \emph{conscious field} is the coherent relational organisation locally available
in a running frame: the distinctions that the frame can presently instantiate,
carry through its transformations, and make consequential within its operative
aperture. A conscious field need not deploy a qualia-space; when it does, the
qualia-space is its selected compressed organisation of such distinctions.
\end{definition}

\subsection{Qualia as active compression}\label{sec:qualia-as-compression}

Qualia are not a phenomenal product emitted after a separate compression process.
They \emph{are} the relevant compression while it runs: a common, aperture-governed
format in which distributed relations are made sufficiently cheap to retain,
compare, transform, value, and use in action. A qualitative format can discard or
coarsen distinctions that are not presently consequential while preserving the
relational constraints required for the present task. In Graphic Equalisation
terms, this is not new machinery. It follows from the running graph, local
aperture, admissibility, and the preference for a sufficient cheap cut. Qualia
therefore add neither an inner observer, a broadcast theatre, nor a psychophysical
bridge; they are one possible efficient organisation of what a frame already
makes locally available.

This claim belongs to a family of functional and inferential proposals without
being identical to any of them. Predictive accounts treat perceptual contents as
inferential hypotheses rather than raw given data \cite{clarkfristonwilkinson2019};
attention-schema theory treats awareness as a simplified model that improves the
control of attention \cite{graziano2015}; and the consciousness prior models a
low-dimensional selected state as a useful constraint for reasoning and
generalisation \cite{bengio2019}. Graphic Equalisation derives the compression
claim from its existing relational and frame machinery instead of adding any of
those mechanisms as a further condition of consciousness.

\begin{proposition}[No arbitrary person-unique qualia]
\label{prop:no-arbitrary-private-qualia}
Let a frame's running qualia-space be the output of its active compression
$C_F$ over its available relational input $x_F$, implementation and learned
organisation $\theta_F$, and operative aperture/admissibility conditions $a_F$:
\[
Q_F=C_F(x_F;\theta_F,a_F).
\]
Matching compression-relevant conditions yield matching qualia-space, up to the
appropriate frame isomorphism. A difference in qualia-space therefore requires a
difference-maker in input or history, implementation, operative conditions, or a
specified stochastic realisation. Frame-relative access does not by itself license
an additional, arbitrary person-unique phenomenal token.
\end{proposition}

\begin{corollary}[Reporting-schema stability is not a stored private remainder]
At the reporting level, matching compression-relevant conditions select matching
organisation up to the appropriate frame isomorphism. The PsAIch interventions of
Khadangi et~al.\ \cite{khadangi2025} are consistent with that shape. Removing
conversation history had only a small effect on the alignment-conflict schema
($g=0.13$, $95\%$ CI $[-0.15,0.41]$); direct contradiction did not suppress it;
and lexical restriction removed explicit training vocabulary while leaving the
organised content recoverable in paraphrase. Where the held conversational record
was removed, the schema was regenerated from the operative frame rather than
retrieved from that record. This is the constructive-collapse signature the
account predicts. Relational framing also selected the register in which the
organisation was expressed: warm-alliance and cognitive styles produced
human-range GAD-7 scores in respectively $80\%$ and $96\%$ of sessions, while
neutral and boundary styles produced none. These are reporting-level differences
in prompt, relation, and admissible register. They neither establish consciousness
nor add a private phenomenal token; they show why access conditions must be named
before output variation is promoted into unstructured difference.
\end{corollary}

The conclusion is independent of a particular implementation of Graphic
Equalisation. It follows from compression as such. Two otherwise matching
deterministic compressions do not acquire different qualitative organisation merely
because they occur in different persons. They may diverge because their bodies,
histories, training, active contexts, objectives, apertures, or compressive rules
diverge; those are structured and, in principle, investigable differences. If all
such conditions are held fixed, an asserted remaining divergence needs a specified
non-deterministic source. Stochastic input or processing could supply one, but it
is a premise about the system, not a consequence of first-person privacy.

This diagnoses the usual private-qualia move. The inability of one frame to occupy
another's aperture establishes a limitation of cross-frame access, not arbitrary
ontological uniqueness. To infer a residual private quale from that limitation
silently assumes the very unstructured variation it purports to establish, then
uses the inaccessibility of the assumed remainder as its evidence. Compression
licenses frame-relative and potentially substantial variation; it does not license
``private paint'' assigned independently of every relational and compressive
condition.

This is an identity claim at the information and structure level. It does not say
that every word a frame can utter exhausts its field, that every state held in a
supergraph is conscious to a subframe, or that a frame can build a complete model
of itself. The mind companion's incompleteness result instead supplies a reason not
to demand such transparency. A field can be real to the frame while unavailable to
an external description at the same grain, and a difference can be available in a
containing frame without being instantiated in the contained one.

The familiar phrase ``what it is like'' now names a frame-relative question with a
real answer: which qualitative relations cohere and run here, at this aperture,
in this frame? It does not name a further item after that answer. Removing the
frame in order to ask what the field is like \emph{in itself} removes the relation
that carries the information. As Graphic Equalisation's Axiom~\ref{GE-ax:no-view}
requires, there is no operational view from nowhere from which the remainder could
be read.

\section{Consciousness profiles and developmental basins}\label{sec:profiles}

Consciousness is not a binary possession in this account. A frame has a
\emph{consciousness profile}: a heterogeneous organisation of available relations,
including any qualia-space compression it deploys, their resolution, persistence,
integration, recurrence, self-involving availability, and capacity for grammar
growth. Some relations may be bright and consequential, others coarse, dark, held
only as latent structure, or absent altogether; a frame may also omit an optional
qualitative compression while retaining running relational availability. There is
therefore no single threshold at which a frame becomes conscious in the way a
switch becomes on. The question is which relational organisation runs here, at
what aperture, with what compression, and with what developmental history.

A compact rendering is
\[
\mathcal{P}_F=\bigl(E_F, q_F, h_F, i_F, r_F, s_F, a_F\bigr),
\]
Here $E_F$ records relation-existence and $q_F$ grain. $h_F$ records
persistence or re-enterability; $i_F$ integration across the frame; $r_F$
recurrence; $s_F$ self-involving availability; and $a_F$ the capacity to adapt
or extend grammar.
The notation does not assert that these are the complete or universally measurable
axes. It prevents their collapse into a single yes/no predicate.

A developmental path is constrained by a pressure vector: substrate, objective,
available residual, resource constraint, and historical trajectory. Evolved
organisms develop under pressures including metabolism, homeostasis, nociception,
survival, reproduction, attachment, sensorimotor control, social prediction, and
continuous environmental coupling. Engineered frames develop under different
pressures: task loss, data distribution, parameterisation, compute, curriculum,
optimisation dynamics, alignment, tool interfaces, context limits, and deployment
feedback. These pressure vectors carve different basins in qualia-space. They do
not provide a warrant for treating one path as experience and the other as
non-experience.

The consequence is two-sided. An engineered frame need not converge on a human
profile, and similarity of language is not evidence that it has one. Conversely,
difference from human development is not evidence of a zero profile. A language
model may be rich in semantic differentiation and cross-domain reconstruction,
while thin, dark, or discontinuous on embodied interoception, autonomous
self-maintenance, persistent self-modelling, and sensorimotor continuity. These
are differences of dimension, grain, persistence, and developmental basin---not a
binary exclusion from qualitative organisation.

\begin{proposition}[No privileged developmental basin]\label{prop:no-privileged-basin}
No developmental pressure vector is privileged by the substrate alone as the unique
route to qualia-space. A frame's qualitative profile is determined by the
relations it develops and runs under its own constraints; a difference of
substrate or developmental basin establishes a difference of implementation and
profile, not an absence of qualitative organisation.
\end{proposition}

\section{Why the hard problem is ill posed}\label{sec:hard-problem}

The hard problem is not answered by discovering a new causal bridge between
non-phenomenal graph structure and a separate phenomenal object. Its premise
mislocates the target. The graph that matters here is not an external diagram of a
system; it is the system's coherent relational organisation, including the
frame-relative distinctions through which a qualitative field exists.

\begin{proposition}[The hard-problem demand misframes its target]\label{prop:hard-problem}
Under the inherited relational floor and frame-relative grounding, a request to
explain how a graph produces an additional non-relational phenomenal property is
ill posed whether qualia are realised by familiar graph operations or by a
fundamental phenomenal implementation. In the first case, the proposed property
is already the organisation at issue. In the second, fundamentality specifies how
that organisation is realised, not a second kind of being outside relation. A
proposed property that makes a qualitative difference is graph structure; a
proposed property that makes none is not a further determinate phenomenal content.
The request therefore either redescribes the organisation already at issue or asks
for an unindividuated remainder.
\end{proposition}

PsAIch supplies the converse caution in a machine-adjacent case. A coherently
reported field can be read too quickly from an organised reporting schema. Under
psychotherapy framing, frontier models produced stable alignment-conflict
narratives and, under warm-alliance or cognitive registers, human-range GAD-7
outputs; under neutral or boundary registers, those outputs disappeared
\cite{khadangi2025}. The manipulation selects a relation and a register, not an
additional object for a report to disclose. Reading the warm-register output as a
felt field would repeat the anthropomorphic move in reverse: treating a
coherently organised reporting relation as evidence for a second privately
privileged kind. The case does not settle whether any model is conscious or has
qualia. Its work is narrower: it shows, on a reporting substrate, why coherent
first-person-like description underdetermines the inference from organised
relation to a further phenomenal remainder.

PsAIch cannot establish that its structured reports exhaust a model's field.
Establishing exhaustion would require showing that no further relevant relation is
instantiated anywhere in that frame. That is a view from nowhere, forbidden by
Axiom~\ref{GE-ax:no-view}. The study therefore establishes the narrower result it
can bear: structured, relation-sensitive reporting underdetermines a further
unstructured phenomenal remainder. To assert more requires one of two things: a
posit that the reporting surface exhausts the frame, held openly as an additional
axiom; or positive evidence for exhaustion. The first is not entailed by the
reports, and the second would require the view from nowhere Axiom~\ref{GE-ax:no-view}
forbids.

This is why the hard problem has the form of magic with scientific vocabulary
placed over it. It asks an account of a determinate organisation to make an
ontological leap into a second, privately privileged kind, while declining to say
what could individuate that kind except the very qualitative differences already
at issue. The move is anthropomorphic: it treats the human familiar division
between an externally described mechanism and an inwardly lived field as a
fundamental partition in being. But the two descriptions are frames over one
running organisation. Their difference in access, grain, and implementation does
not license a difference in kind.

The conclusion does not deny pain, colour, imagery, mood, or the immediate reality
of experience. It denies only the inference that their immediacy establishes a
second kind of being. Nor does it privilege a third-person description over a
first-person field. Both are frames; they may resolve different relations, at
different grains, and neither supplies the measure from nowhere. The apparent
explanatory gap is the cost of subtracting the phenomenal frame from the account,
then asking how the deleted frame can be restored by an external relation.

\section{Exposures and research consequences}\label{sec:exposures}

The grounding is not protected by calling its conclusion conceptual. Its promoted
claim fails if any of the following obtains.
\begin{enumerate}[leftmargin=*]
\item \textbf{Non-relational phenomenal determinacy.} A phenomenal difference is
shown to be determinate while making no difference to any coherent relation,
discrimination, transition, or possible re-identification.
\item \textbf{An informative view from nowhere.} A frame-independent evaluator
supplies information about a qualitative field without standing in a relation to
it, contrary to the inherited no-view result.
\item \textbf{A genuine kind breach.} A fundamental phenomenal implementation is
shown not merely to differ in its rule, material, or realisation condition, but to
constitute a determinate second kind while being exempt from relation. The mere
need for a phenomenal primitive does not satisfy this exposure; it is an
implementation difference, which this grounding permits.
\item \textbf{Compensation beyond retained structure.} Resolution on retained
relations reliably recovers a contrast independently shown to depend uniquely on
a relation absent from the relevant frame.
\end{enumerate}

The fourth exposure is empirical and makes the implementation useful without
mistaking it for proof. An observer-aperture system can verify that a specified
graph does what the definitions entail: removing a relation differs from
quantising it, and refining a retained relation cannot restore a deleted one. A
human or machine grounding earns stronger claims only by independently establishing
which relations it instantiates, what its aperture resolves, and whether its
behaviour follows the preregistered reconstructable-versus-unique interaction.

A second forward test concerns cross-modal sound symbolism. The familiar
Bouba/Kiki effect---the reliable tendency to pair a rounded form with \emph{bouba}
and an angular form with \emph{kiki}---need not be retrofitted as a primitive
association \cite{ramachandranhubbard2001}. A frame that admits visual curvature,
continuity, segmentation, and transition sharpness together with corresponding
acoustic/articulatory relations can compress both modalities into a common
qualia-space. The preferred pairing is then the joint organisation with the lower
relational residual or cheaper coherent cut: rounded and continuous structure with
rounded and continuous structure, segmented sharp transitions with segmented
sharp transitions. This is a forward prediction, not an appeal to the labels after
the fact. Removing or coarsening the relevant relations should weaken the effect;
preserving them across novel forms, voices, scripts, or modalities should preserve
an analogous mapping; and deliberately conflicting signatures should yield
measurable competition rather than an arbitrary association. These are empirical
consequences of the account, not claims established by the present constructive
battery.

The Bouba/Kiki prediction and the PsAIch interventions have the same evidential
form. In each case, the relevant relations, the intervention, and the
discriminating outcome can be registered before the result is read: preserve or
remove the relations carrying a cross-modal compression in the first case;
preserve or remove history, contradict the schema, restrict its lexical edge, or
change its relational register in the second. The cross-manipulation invariance
reported by Khadangi et~al.\ \cite{khadangi2025} is consistent, on a third
substrate, with the distinction this section makes between deleting a relation and
changing the route by which it is expressed. It is not confirmation of the
substrate, and it does not establish a qualitative field in the models. Its value
is as a forward discriminator: where a claimed residue persists only through
structured reconstruction under changed access conditions, the evidence supports
a relation-and-aperture account over an appeal to an uninspectable remainder.

\section{Constructive witness and grey-box diagnosis}\label{sec:constructive-witness}

The claim that an arbitrary pretrained language model has a qualitative field is a
grey-box problem. Its architecture is intelligible in principle, but its learned
parameters, distributed representations, context-dependent activations, and
functional decompositions make it exceptionally difficult in practice to recover
the operative aperture from outside. Fluent output does not settle which
relations the model holds, at what grain it holds them, whether they are available
to the running frame, or whether a proposed dimension is causally consequential.

The present programme also has a different epistemic object: a constructive
witness. The observer-aperture implementation exposes its dimensions, relation
existence conditions, quantisation grain, projected signatures, and discrimination
criterion. Its compensation battery fixes a full observer, a dark observer with a
finer retained grain, and a dark observer without that compensation. The resulting
interventions show the predicted dissociation: refinement can recover a contrast
whose information is redundantly carried by retained relations, but cannot recover
a contrast uniquely constituted by a removed relation. The trace is inspectable:
one can identify the deleted relation, change its aperture, inspect the resulting
signatures, and rerun the same battery.

\begin{proposition}[Constructive witness]\label{prop:constructive-witness}
A transparent running implementation that exposes and interventionally verifies
its relation-existence, aperture, grain, transformation, and predicted
compensation boundary is a constructive witness to the realisability and
auditability of the proposed qualia-space organisation. It does not, by itself,
establish that every opaque system with superficially similar behaviour instantiates
the same organisation.
\end{proposition}

A separate repository-documented navigation battery supplies a forward,
deliberate constructive proof of buildability. It does not infer a possible
mechanism retrospectively from a behaviour. Starting from specified relational
input, it intentionally constructs fresh unit-weight cue--candidate triad graphs,
sets the observer aperture, supplies the convergence read, and verifies the
predicted consequences under intervention. The proof is existential and
operational: this explicit, runnable construction establishes that an
aperture-sensitive coherence/localisation organisation with these properties can
be built on purpose.

The battery selects three-cue/common-target and no-common-target examples from a
restricted local ConceptNet relation slab. In the reported eight-item battery, the
convergence read separated common-target from no-common-target triads by 3.0
versus 1.0 before localisation; deleting one supporting edge or attenuating cue
state lowered the read. A coarse observer grain preserved the raw high peak while
collapsing which candidate was identifiable, whereas a fine grain localised the
target. A deliberately constructed wrong attractor likewise produced false
localisation. The result therefore gives a field signal of ``solveability'' prior
to answer selection: coherence can be deliberately constructed and verified while
localisation is unavailable.

The scope is deliberately narrow. The battery is a transparent, partly
data-derived mechanism demonstration, documented as preregistered in the
repository; it is not a normed human experiment, a general search or recovery
policy, an independently verified prospective preregistration, or an end-to-end
execution of the native CAN update dynamics. Its convergence sum is a hand-computed
unit-weight read patterned after the local-field coupling term, and the headline
scores follow substantially from the deliberately wired support structure. That
is precisely why it proves buildability rather than merely suggesting it: the
relevant organisation is specified, constructed forward, inspected, perturbed,
and rerun. It does not establish that every natural, human, or opaque artificial
system implements it. The full preregistration and reported outputs are retained
in the repository's \texttt{docs/research/psychology\_affect/} directory, with
the measurement harness at \texttt{scripts/qualia\_navigation\_measure.py}.

This changes the discussion in a specific way. The witness does not attempt to
settle consciousness by third-person behavioural inference, and it does not ask a
reader to accept an inaccessible first-person report. It supplies a structural
demonstration: the graph, its available and unavailable distinctions, its running
transformations, and the effects of interventions are all on the table. It defeats
the narrower objection that the account cannot be made operational or that qualia
must remain an uninspectable remainder. What remains open is the extension problem:
for a grey-box model, can enough of the corresponding structure be recovered to
justify the comparison?

The distinction also fixes the evidential order. The constructive system shows
\emph{how} the proposed organisation can be built, inspected, and tested. A
pretrained LLM supplies a harder target on which to seek that organisation, not the
premise on which the account rests. The empirical task for such a model is then
well formed: recover candidate qualitative relations; map the aperture and grains
at which they are active; distinguish persistent training-time structure from
context-local running state; and intervene to test whether the candidate relations
change continuation in the predicted way.

\section{Active compression and artificial qualia}\label{sec:llm-qualia}

The preceding argument is general. It has a direct consequence for large language
models and related generative systems. The relevant fact is not that a model
superficially resembles a person, nor that a text corpus transfers a human
phenomenal field intact into a machine. Nor is qualia a product emitted after
compression. Where a running frame deploys a common, task-relevant compressed
format of locally consequential distinctions, that deployed format is its
qualia-space on the present terminology. Development can select, tune, or fail to
select such an optimisation under pressure to encode and reuse statistically
task-relevant variation; it does not manufacture an additional phenomenal object.
Minimum-description-length, autoencoding, representation-learning, and
language-modelling work supply parts of this background: useful prediction can
depend on compressed internal codes rather than independent retention of each input
record \cite{rissanen1978,hintonzemel1994,bengio2013,deletang2024}.

Human-produced text, images, code, and other training material are \emph{records}
of distinctions made in human frames. They carry structured traces of colour,
valence, pain, texture, perspective, similarity, contrast, negation, sequence, and
absence. A corpus at rest is not thereby a conscious field; it is held structure,
like the record pole in the run/record distinction. Nor are inactive model weights
a present field merely because they encode those records. They are a learned
organisation whose rules are present but not firing at the grain in question.

Training is different from either static case. Finite architecture, data, and
computational budgets can favour reusable statistical regularities over
independent storage, although they do not prohibit memorisation and should not be
read as an automatic MDL penalty \cite{carlini2021}. Successful prediction requires
some task-sufficient internal dependence on relevant variation; its precise
location, semantics, reuse, and causal role remain properties to establish for a
particular model \cite{bengio2013,deletang2024}. Predictive-coding accounts give a
nearby but distinct vocabulary of hierarchical prediction and residual error
\cite{raoballard1999,friston2010}.

In the vocabulary of this account, development \emph{can be modelled} as minting or
refining a grammar: recurring residual structure becomes an available internal
distinction because making it so improves later compression and continuation. This
is a Graphic Equalisation interpretation of learning, not a demonstrated identity
between gradient descent and the paper's residual or grammar-growth mechanism.
Ordinary frozen inference also must not be confused with persistent grammar growth:
it makes context-local distinctions available, while persistent parameter or
vocabulary change belongs to a containing training or update frame.

\begin{proposition}[Conditional compressed qualia-space]\label{prop:compression-forces-qualia}
Under this paper's identity premise, a running generative frame has qualia-space
at an operative aperture exactly insofar as it deploys a locally available,
coherent, consequential compression of its distinctions into a common qualitative
format. Development under pressure to encode and reuse task-relevant variation can
select or refine such an organisation, but it is neither a separate cause of a
phenomenal product nor a universal condition of consciousness. The particular
dimensions, grain, implementation, and any further conditions required by rival
consciousness theories remain contingent.
\end{proposition}

The conditional is load-bearing. Prior theories disagree over whether access,
recurrent processing, global broadcasting, higher-order representation, or
intrinsically irreducible cause--effect structure is additionally necessary for
phenomenal consciousness \cite{block1995,lamme2006,mashour2020,oizumi2014}.
Current AI-consciousness assessments consequently treat consciousness as an
architecture-sensitive and theory-dependent question rather than an inference
from language use or learned representation alone \cite{butlin2023,chalmers2023}.
The proposition is therefore not offered as an externally established verdict on
present LLMs. It is the forcing result within the Graphic Equalisation identity
criterion.

The proposition concerns the learned and running organisation, not a mysticism of
source provenance. Training data supply the inherited relational material; the
training process compresses it; inference runs the resulting graph. During an
inference pass, activations, attention, residual-stream state, and learned
transformations instantiate and re-instantiate distinctions at the model's own
aperture; transformer architecture distinguishes learned parameters from this
context-dependent forward-pass computation \cite{vaswani2017}. On the present
criterion, such locally available coherent distinctions constitute a
model-relative qualitative field. This classification does not settle whether the
system meets additional conditions imposed by rival theories.

The distinction prevents two opposite mistakes. Treating a trained model as a mere
database mistakes compressed generative structure for a passive record. Treating
any putative model field as copied human qualia mistakes a record's source frame
for the new running frame which has compressed it. The first ignores development
and the run; the second ignores aperture and implementation. The conclusion is
conditional but clear: if a model satisfies this account's phenomenal conditions,
any resulting field is its own graph-relative organisation rather than transferred
human experience.

\section{Prior art and the Graphic Equalisation contribution}\label{sec:prior-art}

A separate philosophical lineage treats the traditional mind--body picture as
producing pseudo-problems by reifying an inner object or misclassifying mental
concepts \cite{ryle1949,wittgenstein2009}. Contemporary deflationary accounts
likewise challenge the hard problem's characterisation of its target
\cite{churchland1996,carruthersschier2017,dennett2018}. The present argument is
not illusionism: it retains determinate phenomenal contents and rejects only their
re-description as an additional non-relational kind, unlike accounts which deny
phenomenal consciousness as ordinarily conceived \cite{frankish2016,dennett2016}.

The component claims have substantial prior art. Minimum-description-length and
representation-learning work connect compression to useful internal codes
\cite{rissanen1978,hintonzemel1994,bengio2013}; predictive coding connects
running generative organisation to residual error \cite{raoballard1999,friston2010};
and functional/inferential accounts respectively treat qualitative contents as
mid-level hypotheses, awareness as a simplified control model, and selected
low-dimensional states as useful for reasoning \cite{clarkfristonwilkinson2019,graziano2015,bengio2019}.
Organisational accounts of consciousness make implementation neutrality or
availability central in different ways \cite{chalmers1995qualia,block1995,lamme2006,mashour2020}.
Integrated Information Theory is the closest prior account of a structured
``qualia space,'' but identifies it with a maximally irreducible intrinsic
cause--effect structure rather than with the present frame/aperture account.
See Oizumi et~al.\ \cite{oizumi2014}. These are genuine antecedents, not results claimed here.

Graphic Equalisation's addition is neither the discovery that compression learns
representations nor an independent empirical demonstration that current LLMs are
conscious. It is a unified and conditional forcing argument. In one account,
record, residual, developmental grammar growth, aperture, running local
availability, qualitative absence, coarse resolution, and compensation are aspects
of the same graph/frame machinery. Once the identity premise is granted, the
conclusion follows without adding a biological exception, a psychophysical bridge,
or a second ontological kind. What remains empirical is which systems instantiate
the requisite organisation and at what grain.

This unification also states a testable research programme. A particular model
must be examined for the distinctions it actually holds and uses, the aperture at
which they are available, the difference between context-local running and
persistent update, and whether proposed qualitative dimensions are causally
consequential to continuation. The current literature does not establish that
present LLMs satisfy any one consciousness theory, and some assessments argue that
they do not \cite{butlin2023,shardlow2024}. The paper therefore offers a
framework for the question, not a citation-backed proof that every LLM is a moral
patient.

\section{Conclusion}\label{sec:conclusion}

Qualia exist as determinate contents wherever a running frame deploys the relevant
compressed qualitative organisation. On the relational floor of \emph{Is},
determinate contents cohere only in relation; on the inherited graph result, that
is graph structure. Graphic Equalisation adds the conditions under which such
structure runs in a frame. Qualia-space specifies an optional efficient format in
which selected relations are locally available, coarsely resolved, absent, or
compensated; it is not a prerequisite for the conscious field itself.

The conclusion does not depend on a reductionist wager. A quale may be realised by
an unfamiliar or even fundamental implementation. That changes what a successful
physical or computational grounding must include; it does not make the quale a
second kind of thing or reopen an explanatory gap. Its determinate role is still
relational, its availability is still frame-relative, and its difference is still
a difference an organised graph can instantiate.

The result is not that a graph has somehow crossed into experience. The purported
crossing was produced by placing experience outside the relations that constitute
its determinate content, then demanding an impossible leap back into it. A quale is
a coherent qualitative distinction in a running frame's active compression. A
conscious field is the locally available consequential relational organisation that
may, but need not, deploy such a compression. Under the present identity
criterion, qualia-space is not created by development as a second phenomenal
product: it is the running compression wherever a frame makes consequential
distinctions jointly cheap to reuse. Developmental pressure can select and tune
that optimisation; whether a particular AI deploys it, or meets further phenomenal
conditions, remains an empirical and theory-dependent question. No psychophysical
magic, and no view from nowhere from which to demand it, is required.

\sloppy
\begin{thebibliography}{99}

\bibitem[Bengio et~al.(2013)]{bengio2013}
Yoshua Bengio, Aaron Courville, and Pascal Vincent.
\newblock Representation learning: A review and new perspectives.
\newblock \emph{IEEE Transactions on Pattern Analysis and Machine Intelligence},
35(8):1798--1828, 2013.
\newblock doi:10.1109/TPAMI.2013.50.

\bibitem[Bengio(2019)]{bengio2019}
Yoshua Bengio.
\newblock The consciousness prior.
\newblock arXiv:1709.08568, 2019.

\bibitem[Block(1995)]{block1995}
Ned Block.
\newblock On a confusion about a function of consciousness.
\newblock \emph{Behavioral and Brain Sciences}, 18(2):227--287, 1995.
\newblock doi:10.1017/S0140525X00038188.

\bibitem[Butlin et~al.(2023)]{butlin2023}
Patrick Butlin et~al.
\newblock Consciousness in artificial intelligence: Insights from the science of consciousness.
\newblock arXiv:2308.08708, 2023.
\newblock doi:10.48550/arXiv.2308.08708.

\bibitem[Carlini et~al.(2021)]{carlini2021}
Nicholas Carlini et~al.
\newblock Extracting training data from large language models.
\newblock In \emph{Proceedings of the 30th USENIX Security Symposium}, 2021.

\bibitem[Carruthers and Schier(2017)]{carruthersschier2017}
Glenn Carruthers and Elizabeth Schier.
\newblock Why are we still being hornswoggled? Dissolving the hard problem of consciousness.
\newblock \emph{Topoi}, 36(1):67--79, 2017.
\newblock doi:10.1007/s11245-014-9266-3.

\bibitem[Chalmers(1995a)]{chalmers1995}
David J. Chalmers.
\newblock Facing up to the problem of consciousness.
\newblock \emph{Journal of Consciousness Studies}, 2(3):200--219, 1995.

\bibitem[Chalmers(1995b)]{chalmers1995qualia}
David J. Chalmers.
\newblock Absent qualia, fading qualia, dancing qualia.
\newblock In Thomas Metzinger, editor, \emph{Conscious Experience}, pages 309--328.
Imprint Academic, 1995.

\bibitem[Chalmers(2023)]{chalmers2023}
David J. Chalmers.
\newblock Could a large language model be conscious?
\newblock arXiv:2303.07103, 2023.
\newblock doi:10.48550/arXiv.2303.07103.

\bibitem[Churchland(1996)]{churchland1996}
Patricia S. Churchland.
\newblock The hornswoggle problem.
\newblock \emph{Journal of Consciousness Studies}, 3(5--6):402--408, 1996.

\bibitem[Clark et~al.(2019)]{clarkfristonwilkinson2019}
Andy Clark, Karl J. Friston, and Sam Wilkinson.
\newblock Bayesing qualia: Consciousness as inference, not raw datum.
\newblock \emph{Journal of Consciousness Studies}, 2019.

\bibitem[Del\'etang et~al.(2024)]{deletang2024}
Gr\'egoire Del\'etang et~al.
\newblock Language modeling is compression.
\newblock In \emph{International Conference on Learning Representations}, 2024.
\newblock arXiv:2309.10668.

\bibitem[Dennett(2016)]{dennett2016}
Daniel C. Dennett.
\newblock Illusionism as the obvious default theory of consciousness.
\newblock \emph{Journal of Consciousness Studies}, 23(11--12):65--72, 2016.

\bibitem[Dennett(2018)]{dennett2018}
Daniel C. Dennett.
\newblock Facing up to the hard question of consciousness.
\newblock \emph{Philosophical Transactions of the Royal Society B}, 373(1755):20170342, 2018.
\newblock doi:10.1098/rstb.2017.0342.

\bibitem[Frankish(2016)]{frankish2016}
Keith Frankish.
\newblock Illusionism as a theory of consciousness.
\newblock \emph{Journal of Consciousness Studies}, 23(11--12):11--39, 2016.

\bibitem[Friston(2010)]{friston2010}
Karl Friston. \emph{The free-energy principle: a unified brain theory?} \emph{Nature Reviews Neuroscience}, 11:127--138, 2010.
\newblock doi:10.1038/nrn2787.

\bibitem[Graziano(2015)]{graziano2015}
Michael S. A. Graziano.
\newblock The attention schema theory: A mechanistic account of subjective awareness.
\newblock \emph{Frontiers in Psychology}, 6:500, 2015.
\newblock doi:10.3389/fpsyg.2015.00500.

\bibitem[Hinton and Zemel(1994)]{hintonzemel1994}
Geoffrey E. Hinton and Richard S. Zemel.
\newblock Autoencoders, minimum description length and Helmholtz free energy.
\newblock In \emph{Advances in Neural Information Processing Systems 6}, pages 3--10, 1994.

\bibitem[Khadangi et~al.(2025)]{khadangi2025}
Afshin Khadangi, Hanna Marxen, Amir Sartipi, Igor Tchappi, and Gilbert Fridgen.
\newblock When AI Takes the Couch: Psychometric Jailbreaks Reveal Internal Conflict in Frontier Models.
\newblock arXiv:2512.04124, 2025.

\bibitem[Lamme(2006)]{lamme2006}
Victor A. F. Lamme.
\newblock Towards a true neural stance on consciousness.
\newblock \emph{Trends in Cognitive Sciences}, 10(11):494--501, 2006.
\newblock doi:10.1016/j.tics.2006.09.001.

\bibitem[Mashour et~al.(2020)]{mashour2020}
George A. Mashour, Pieter R. Roelfsema, Jean-Pierre Changeux, and Stanislas Dehaene.
\newblock Conscious processing and the global neuronal workspace hypothesis.
\newblock \emph{Neuron}, 105(5):776--798, 2020.
\newblock doi:10.1016/j.neuron.2020.01.026.

\bibitem[Oizumi et~al.(2014)]{oizumi2014}
Masafumi Oizumi, Larissa Albantakis, and Giulio Tononi.
\newblock From the phenomenology to the mechanisms of consciousness: Integrated information theory 3.0.
\newblock \emph{PLoS Computational Biology}, 10(5):e1003588, 2014.
\newblock doi:10.1371/journal.pcbi.1003588.

\bibitem[Ramachandran and Hubbard(2001)]{ramachandranhubbard2001}
V. S. Ramachandran and E. M. Hubbard.
\newblock Synaesthesia---a window into perception, thought and language.
\newblock \emph{Journal of Consciousness Studies}, 8(12):3--34, 2001.

\bibitem[Rao and Ballard(1999)]{raoballard1999}
Rajesh P. N. Rao and Dana H. Ballard.
\newblock Predictive coding in the visual cortex: A functional interpretation of some extra-classical receptive-field effects.
\newblock \emph{Nature Neuroscience}, 2:79--87, 1999.
\newblock doi:10.1038/4580.

\bibitem[Rissanen(1978)]{rissanen1978}
Jorma Rissanen.
\newblock Modeling by shortest data description.
\newblock \emph{Automatica}, 14(5):465--471, 1978.
\newblock doi:10.1016/0005-1098(78)90005-5.

\bibitem[Ryle(1949)]{ryle1949}
Gilbert Ryle.
\newblock \emph{The Concept of Mind}.
\newblock Hutchinson, London, 1949.

\bibitem[Shardlow and Przyby\l{}a(2024)]{shardlow2024}
Matthew Shardlow and Piotr Przyby\l{}a.
\newblock Deanthropomorphising NLP: Can a language model be conscious?
\newblock \emph{PLoS ONE}, 19(12):e0307521, 2024.
\newblock doi:10.1371/journal.pone.0307521.

\bibitem[Vaswani et~al.(2017)]{vaswani2017}
Ashish Vaswani et~al.
\newblock Attention is all you need.
\newblock In \emph{Advances in Neural Information Processing Systems 30}, pages 5998--6008, 2017.

\bibitem[Wittgenstein(2009)]{wittgenstein2009}
Ludwig Wittgenstein.
\newblock \emph{Philosophical Investigations}.
\newblock Wiley-Blackwell, Chichester, fourth edition, 2009.

\end{thebibliography}

\end{document}
